\chapter{Додатковий нагрів і генерація струму}\label{ch:heating}

Плазму догрівають нейтральними пучками (NBI), хвилями електронно-циклотронного (ECH), іонно-циклотронного (ICRH) і нижньогібридного (LH) діапазонів та α-частинками. Ті самі засоби генерують струм і формують профіль $q$ (розд.~\ref{ch:stability}).

\section{Омічний нагрів і його межа}\label{sec:10-ohm}

\begin{outofcorpus}[опір Спітцера і неокласичний опір]
Паралельний опір Спітцера $\eta_{\mathrm{Sp}}\approx\num{1,65e-9}\,\Zeff\lnL/T_e^{3/2}$~Ом$\cdot$м ($T_e$ у кеВ) спадає разом із частотою зіткнень $\nu_{ei}\propto T_e^{-3/2}$. Захоплені електрони струму не несуть, тому в банановому режимі $\eta_{\mathrm{neo}}\approx\eta_{\mathrm{Sp}}/(1-\sqrt{\varepsilon})^2$. Потужність $P_\Omega=\eta j^2\propto T_e^{-3/2}$, а втрати ростуть з $T$: омічна фаза насичується на $T_e\sim1$--3~кеВ. Час скінування $\tau_R=\mu_0a^2/\eta\sim\SI{4000}{с}$ при \SI{10}{кеВ}, $a=\SI{2}{м}$ (п.~\ref{sec:05-15d}).
\end{outofcorpus}

Звідси «пам'ять» профілю струму на сотні секунд у TASK/TR~\cite[с.~4, 10]{Fukuyama1992}, рівн.~\eqref{eq:05-Bp}. На WT-3 ($\Ip=\SI{80}{кА}$) ECH підняв $T_e(0)$ з 510 до 670~еВ, а $\Vloop$ упала з 1,5 до 0,8~В~\cite[с.~4]{Hanada1990} (дані для LFS, $q_L=4{,}3$; задача~1).

\section{Інжекція нейтральних пучків}\label{sec:10-nbi}

Нейтрали пучка іонізуються (зокрема багатоступенево): уздовж променя $\dd N/\dd l=-n_e\sigma_bN$, де $\sigma_b(E_b,n_e,n_i,\Zeff)$~--- переріз гальмування пучка. Спад $N$ дає джерело швидких іонів і потужність $P_b$~\cite[с.~6, (15)]{Fukuyama1992}. Швидкі іони сповільнюються на своїй магнітній поверхні, і їхній розподіл має вигляд~\cite[с.~4, (7)--(9)]{Fukuyama1992}
\begin{equation}\label{eq:10-fsd}
  f(v)=\frac{S_b\tau_{\mathrm s}}{4\pi}\,\frac{\Theta(v_b-v)}{v^3+v_{\mathrm{cr}}^3},\qquad
  \tau_{\mathrm s}=\frac{6\pi\sqrt{2\pi}\,\varepsilon_0^2m_bT_e^{3/2}}{n_eZ_b^2m_e^{1/2}e^4\lnL},\qquad
  v_{\mathrm{cr}}^3=\sum_i3\sqrt{\frac{\pi}{2}}\,\frac{n_iZ_i^2}{n_e}\,\frac{m_e}{m_i}\Big(\frac{T_e}{m_e}\Big)^{3/2},
\end{equation}
де $S_b$~--- темп народження, $\tau_{\mathrm s}$~--- час гальмування на електронах; при $v=v_{\mathrm{cr}}$ (у джерелі $v_c$) гальмування на іонах таке саме.

\begin{derivation}[критична енергія і розподіл потужності]
Сповільнення: $\dot v=-\tau_{\mathrm s}^{-1}(v+v_{\mathrm{cr}}^3/v^2)$, де перший доданок дають електрони, другий~--- іони. З \eqref{eq:10-fsd} $E_{\mathrm{cr}}=\tfrac12m_bv_{\mathrm{cr}}^2\approx14{,}8\,A_bT_e\big[\sum_in_iZ_i^2/(n_eA_i)\big]^{2/3}$ (не плутати з полем $\Ec$). Для D у D це $18{,}6\,T_e$, у D--T 50:50~--- $16{,}5\,T_e$ для D і $32{,}8\,T_e$ для α. Поділивши іонну частину $\dd E/\dd t$ на повну і проінтегрувавши по $\dd v$, дістанемо частку енергії, що йде іонам:
\begin{equation}\label{eq:10-H}
  h_i(y)=\frac{2}{y^2}\int_0^y\frac{x\,\dd x}{x^3+1},\qquad y=\frac{v_b}{v_{\mathrm{cr}}};\qquad h_i\to1\ (y\ll1),\quad h_i\to0\ (y\gg1).
\end{equation}
Це функція $H$ з~\cite[с.~5, (11)--(12)]{Fukuyama1992}. TASK/TR замість $f(v,t)$ розв'язує $\dd W_b/\dd t=P_b-W_b/\tau_b$ з $\tau_b=\tfrac12\tau_{\mathrm s}[1-h_i(y)]$, що дає правильну стаціонарну енергію $W_b$. Електрони отримують $(W_b/\tau_b)(1-h_i)$, іони~--- $(W_b/\tau_b)h_i$~\cite[с.~4--5, (10)--(14)]{Fukuyama1992}.
\end{derivation}

\textbf{Генерація струму.} Струм швидких іонів $j_i$ частково екранують електрони: $j_t=j_i(1-Z_b/\Zeff)$. Захоплені електрони послаблюють екранування~\cite[с.~6, (16)--(17)]{Fukuyama1992}, і ефективність генерації становить $j/P_b=(2eZ_b\tau_{\mathrm s}/m_bv_{\mathrm{cr}})\allowbreak[1-(Z_b/\Zeff)(1-G_{\mathrm t})]F_{\mathrm{FP}}$ з поправками Фоккера--Планка $G_{\mathrm t}(\Zeff,\varepsilon)$, $F_{\mathrm{FP}}$.

\textbf{Повна генерація струму в ITER-класі (TASK/TR).} $R=\SI{6}{м}$, $a=\SI{2}{м}$, $\Bzero=\SI{4,85}{Тл}$, $\Ip=\SI{22}{МА}$, $E_b=\SI{1}{МеВ}$~\cite[с.~7]{Fukuyama1992}. При \SI{150}{МВт} і $\fsa{n_e}\approx\SI{0,7e20}{м^{-3}}$ поглинена потужність і струм пучка максимальні біля $r\approx\SI{0,35}{м}$ (радіус дотику пучка), бутстреп-струм тече на периферії. Омічний струм зникає при \SI{130}{МВт} ($\fsa{n_e}\approx\SI{0,7e20}{м^{-3}}$) і при \SI{150}{МВт} (\SI{1,2e20}{м^{-3}})~\cite[с.~10, 12]{Fukuyama1992}. Момент пучка~--- розд.~\ref{ch:flows}. У тренажері Fenix (ASTRA) NBI і ECRH рахують зовнішні бібліотеки RABBIT і TORBEAM~\cite[с.~5, рис.~3]{Fable2022}.

\section{Електронно-циклотронний нагрів і генерація струму}\label{sec:10-ech}

\begin{outofcorpus}[резонанс, моди, відсічки]
Резонанс на гармоніці $l$: $\omega-l\Omc{e}/\gamma=\kpar v_\parallel$, $\Omc{e}/2\pi\approx\SI{28}{ГГц}\cdot B[\mathrm{Тл}]$. Оскільки $B\propto1/R$, шар поглинання~--- вертикальна смуга $R_{\mathrm{res}}=lR_0\Bzero e/(m_e\omega)$, яку зсувають вибором $\Bzero$ або $\omega$. О-мода ($\vect{E}\parallel\Bvec$) відсікається при $\omega=\omega_{pe}$, тобто при $N_{\mathrm{cr}}=\varepsilon_0m_e\omega^2/e^2$. X-мода з боку слабкого поля на другій гармоніці доходить до шару, поки не впирається в R-відсічку: $\omega_{pe}^2<\omega(\omega-\Omc{e})$. Поглинання локалізоване в мм--см. ECCD виникає, коли гріють електрони з одним знаком $v_\parallel$: гарячі рідше зіштовхуються (механізм Фіша--Бузера), а захоплення нагрітих дає протилежний внесок (Окава).
\end{outofcorpus}

\textbf{WT-3} ($R=\SI{65}{см}$, $a=\SI{20}{см}$, $\Bzero\le\SI{1,75}{Тл}$). Гіротрон \SI{56}{ГГц} ($P_{\mathrm{ECH}}\le\SI{200}{кВт}$) вводить сфокусовану X-моду з боку слабкого поля під кутом $88^\circ$ до $\Btor$; поглинання на другій гармоніці, шар зсувають зміною $\Bzero$~\cite[с.~3]{Hanada1990}. При $n_e\lesssim\SI{5e18}{м^{-3}}$ (переведено з СГС) з'являються гарячі електрони ($\Teh\approx\SI{60}{кеВ}$, до \SI{300}{кеВ}; пилка~--- розд.~\ref{ch:stability})~\cite[с.~5]{Hanada1990}. Для кондиціонування стінок на Ураган-2М до тліючого розряду додають магнетрон \SI{2,45}{ГГц} (до \SI{0,6}{кВт})~\cite[с.~35]{Martsenyuk2025}; для нього $N_{\mathrm{cr}}\approx\SI{7,45e16}{м^{-3}}$~\cite[с.~112]{Martsenyuk2025}, ЕЦ-резонанс~--- при \SI{0,0875}{Тл} (НВЧ-кондиціонування, п.~\ref{sec:08-gdc}).

\section{Іонно-циклотронний нагрів і швидкі магнітозвукові хвилі}\label{sec:10-icrh}

\begin{outofcorpus}[ICRH]
Антена збуджує швидку магнітозвукову хвилю ($\omega\approx k\vA$), яка проходить у ядро. На основній гармоніці однокомпонентної плазми ліва поляризація при $\omega=\Omc{i}$ зникає, тому гріють меншину (кілька відсотків H чи $^3$He у D) або другу гармоніку. Біля шару меншини лежить іон-іонний гібридний резонанс. Там частина хвилі трансформується (mode conversion) в іонну бернштейнівську хвилю, що гріє електрони. Меншина з енергією $\gg E_{\mathrm{cr}}$ за \eqref{eq:10-H} теж гріє переважно електрони.
\end{outofcorpus}

Умова локального резонансу хвилі з частинкою сорту $s$ на гармоніці $l$~\cite[с.~87--88, (4.1)]{Tyshchenko2021}:
\begin{equation}\label{eq:10-res}
  \omega-l\,\Omc{s}(x)-\omega_{\mathrm d}(r,\theta)=\kpar(r)\,v_\parallel,\qquad \Omc{s}=\bar\Omega_s(1-x/R),\quad x=r\cos\theta,
\end{equation}
де $\omega_{\mathrm d}$~--- частота тороїдального дрейфу (у джерелі $\omega_D$), $\bar\Omega_s$~--- значення на осі. У D і α майже однакова гірочастота (\num{26,7} і \SI{26,9}{МГц} при \SI{3,5}{Тл}~\cite[с.~89]{Tyshchenko2021}). Тому при $l_{\mathrm D}=l_\alpha$ доданки $l\bar\Omega$ скорочуються, і одна швидка магнітозвукова мода (FMM) резонує і з α на периферії, і з дейтронами біля осі. Потрібно лише $v_{\parallel i}^{\mathrm{res}}/v_{Ti}=1{,}5$--2~\cite[с.~89--91]{Tyshchenko2021}. На JET ICRH (0,9 і 2,0~МВт поверх 10~МВт NBI) правив за еталон для α-нагріву~\cite[с.~86, табл.~4.1]{Tyshchenko2021}. При цьому $\tauE$ був найменшим за найбільшої потужності ICRH~\cite[с.~116]{Tyshchenko2021}. Fenix ICRH не моделює: редукованих моделей ICRF немає~\cite[с.~4--5]{Fable2022}.

\section{Нижньогібридний діапазон}\label{sec:10-lh}

\begin{outofcorpus}[НГР і LHCD]
$\varepsilon_\perp=1+\omega_{pe}^2/\Omc{e}^2-\sum_i\omega_{pi}^2/(\omega^2-\Omc{i}^2)$. Резонанс $\varepsilon_\perp=0$ (НГР) настає при $\omega_{\mathrm{LH}}^2\approx\Omc{i}^2+\omega_{pi}^2/(1+\omega_{pe}^2/\Omc{e}^2)$. Повільна хвиля з $\npar>\omega_{pe}/\Omc{e}+\sqrt{1+\omega_{pe}^2/\Omc{e}^2}$ (доступність; при $\omega^2\gg\omega_{pi}^2$, точніше під коренем $-\omega_{pi}^2/\omega^2$) доходить до ядра і затухає за Ландау на електронах з $v_\parallel=c/\npar\approx3v_{Te}$~--- це найефективніша генерація струму (LHCD). Бік, з якого контур обходить сингулярність, задає причинність: дисипація дає $\operatorname{Im}\varepsilon_\perp$ певного знака (правило Ландау).
\end{outofcorpus}

Марчук~\cite[с.~49--50, (2.39)--(2.41)]{Marchuk2025} розв'язує рівняння повільної хвилі для $E_z$ (шар товщини $L$, тестовий профіль $\varepsilon_\perp=x-x_r$):
\begin{equation}\label{eq:10-slow}
  \frac{\dd}{\dd x}\Big(\mathcal F\frac{\dd E_z}{\dd x}\Big)+\mathcal G E_z=0,\qquad \mathcal F=\frac{\npar^2\varepsilon_\perp}{\npar^2-\varepsilon_\perp}\approx\varepsilon_\perp,\qquad \mathcal G=\kpar^2L^2(-\varepsilon_\parallel).
\end{equation}
У НГР $\mathcal F=0$, і метод Гальоркіна з лінійними скінченними елементами не працює. Приклад: D, $B=\SI{2,5}{Тл}$, $n=\SI{e17}{м^{-3}}$, $\omega=\SI{3e8}{с^{-1}}$, $\kpar=\SI{3}{м^{-1}}$, $L=\SI{5}{см}$~\cite[с.~50]{Marchuk2025}. Два обходи сингулярності~\cite[с.~50--53, 58--60]{Marchuk2025}: (1)~\emph{Метод штрафу}: при $|\mathcal F|<F_0$ додають уявну латку, $\mathcal F\to\mathcal F+\mathrm{i}F_{\mathrm{im}}$. Похибка поглиненої потужності сягає 17\,\% для лінійної латки і $<5\,\%$ для параболічної. (2)~\emph{Аналітичне продовження}: коефіцієнти аналітичні, тож шлях зсувають у комплексну площину, $x\to x+\mathrm{i}\ell(x)$, $\ell=x_w\cos^2[\pi(x-x_s)/2x_w]$. Похибка $\lesssim10^{-4}$ при $x_w=0{,}05$--0,25, і косинусний обхід (з неперервною похідною) точніший за параболічний.

\section{α-нагрів і каналювання енергії}\label{sec:10-alpha}

α-частинки несуть $P_\alpha\approx P_{\mathrm{fus}}/5$ (розд.~\ref{ch:trap}) і, оскільки $E_\alpha=\SI{3,5}{МеВ}\gg E_{\mathrm{cr}}$, гріють здебільшого електрони (задача~2). Проте в DT-розрядах JET (DTE1) центральна $T_i$ була вищою, ніж у D-розрядах з ICRH тієї самої потужності~\cite[с.~85]{Tyshchenko2021}. Тищенко пояснює це \emph{доцентровим каналюванням}: FMM, збуджені α на периферії (там народжується $\approx30\,\%$ α, $r/a\ge0{,}5$~\cite[с.~96]{Tyshchenko2021}), переносять частку $\fSC$ потужності $P_\alpha$ до осі й віддають її іонам через \eqref{eq:10-res}~\cite[с.~86--91]{Tyshchenko2021}.

\begin{derivation}[коли каналювання подвоює $T_0$]
Стаціонар $r^{-1}(rn\chi_{\mathrm h}T')'+S_Q^{(c)}+S_Q^{(w)}=0$, $n\chi_{\mathrm h}=\mathrm{const}$, об'ємні джерела $S_Q^{(j)}\propto\exp(-\sigma_jr^2/a^2)$, $\int S_Q^{(w)}\dd^3r=\fSC P_\alpha$, $\int S_Q^{(c)}\dd^3r=(1-\fSC)P_\alpha$. Подвійне інтегрування при $\sigma_c\ll1$ дає~\cite[с.~105--107, (4.45)--(4.51)]{Tyshchenko2021}
\begin{equation}\label{eq:10-T0}
  T_0-T(a)=\frac{a^2P_\alpha}{4n\chi_{\mathrm h}V_p}\Big[1+\fSC\Big(\frac{\mathrm{Ein}(\sigma_w)}{1-\ee^{-\sigma_w}}-1\Big)\Big].
\end{equation}
Тут $\mathrm{Ein}(\sigma)=\int_0^\sigma(1-\ee^{-t})t^{-1}\dd t\approx\ln\sigma+0{,}577$. Дужка дорівнює 2 при $\fSC(\ln\sigma_w-0{,}42)=1$ (наприклад, $\fSC=0{,}53$ при $\sigma_w=10$).
\end{derivation}

Двотемпературна модель JET навіть з $\fSC=0{,}3$ виміряних $T_i$ не відтворює, а структуру FMM на JET не виміряно, тож каналювання~--- гіпотеза, не встановлений механізм (докладно~--- п.~\ref{sec:05-15d})~\cite[с.~89, 115--116]{Tyshchenko2021}.

\subsection*{Контрольні запитання}
\begin{enumerate}
  \item Чому омічний нагрів не доводить до запалювання? Оцініть $\tau_R$ для WT-3 при $T_e=\SI{0,6}{кеВ}$ і порівняйте з тривалістю ECH \SI{30}{мс}~\cite[с.~9]{Hanada1990}.
  \item Навіщо ICRH меншина? Чому з однією FMM резонують α і D, але не тритони~\cite[с.~91]{Tyshchenko2021}?
  \item Оцініть $\omega_{\mathrm{LH}}$ для тесту Марчука. Чому там можна знехтувати $\omega_{pe}^2/\Omc{e}^2$?
\end{enumerate}

\subsection*{Задачі}
\begin{enumerate}
  \item (WT-3~\cite{Hanada1990}.) а)~Знайдіть $B_{\mathrm{res}}$ для \SI{56}{ГГц} (друга гармоніка) і $\Bzero$, за якого шар на осі, на $R_0-\SI{5}{см}$, у межах плазми; де шар при \SI{1,75}{Тл}? б)~Порівняйте $n_e=\SI{5e18}{м^{-3}}$ з відсічками О- і X-моди в шарі. в)~Знайдіть $P_{\mathrm{OH}}=\Vloop\Ip$ до й після ECH і $\Vloop$ за Спітцером ($T_e(0)$: $510\to670$~еВ; дані для LFS, $q_L=4{,}3$; для умов рис.~3~--- лише оцінка). Звідки розбіжність?
  \item (ITER-клас~\cite{Fukuyama1992}: D--T 50:50, $T_e=\SI{10}{кеВ}$, $n_e=\SI{0,7e20}{м^{-3}}$, $\lnL=17$.) Знайдіть $\tau_{\mathrm s}$, $E_{\mathrm{cr}}$, $y$, $h_i$, $\tau_b$ для D 1~МеВ і α 3,5~МеВ та частку 150~МВт NBI, що йде іонам. Порівняйте з D 80~кеВ у D.
\end{enumerate}
